% [AUTHOR REWRITE] interpretation of the estimates
Full attention is rejected in apartments and not in houses. The evidence is the contrast between segments, which does not depend on the normalisation of the default: the perceived cost of silence, identified without it, is \attCstarApt{} kWh per square metre and year in apartments (95\% CI \attCstarLoApt{} to \attCstarHiApt{}) against \attCstarHouse{} (\attCstarLoHouse{} to \attCstarHiHouse{}) in houses ($t=\attDiffCstarT$ for the difference). In apartments it is a small fraction of the \attCsilApt{} kWh a fully attentive buyer would assign to the silent pool, so silent apartments are priced roughly as though energy were not part of the transaction. In houses it is close to the rational cost of the pool, which is why the model of Section~\ref{sec:equilibrium} fits that segment. One attention parameter per segment thus accounts for both the rejection in apartments and the fit in houses, without any change to the seller side; its level, unlike the contrast, depends on the normalisation of the default. By Proposition~\ref{prop:att_identification}(iii), the apartment estimate is best read as an upper bound. House prices are consistent with silence being priced at the rational cost of the pool, which is all Section~\ref{sec:equilibrium} requires; they do not show that house buyers are attentive in any absolute sense.
